\begin{claim}{E.1}{Proved}{What ``one categorical level up'' means}
Let $\mathcal C$ be an $\infty$-category, small in a universe $\V$ larger
than the universe $\U$ of indexing spaces.  Put
\[
 K_0(\mathcal C)=\mathcal C^{\simeq},\qquad
 K_1(\mathcal C)=\Fun(\Delta^1,\mathcal C)^{\simeq}.
\]
For a map $f:X\to Y$ of $\U$-small spaces, restriction
$f^*:\Fun(Y,\mathcal C)\to\Fun(X,\mathcal C)$ is an equivalence if and
only if both canonical maps
\[
 \Map(Y,K_i(\mathcal C))\longrightarrow\Map(X,K_i(\mathcal C)),
 \qquad i=0,1,
\]
are equivalences of spaces.  This is a two-object orthogonality
description, not by itself the constant-coefficient class of a derivator.
If $\mathcal C$ is locally $\U$-small and admits the required
$\U$-small cotensors, equivalence of the
ordinary homotopy categories under the specified restriction functor is
equivalent to the stated $\infty$-categorical condition.  The same
conclusion holds without local $\U$-smallness whenever that
restriction has a Kan-extension adjoint.
\end{claim}
\begin{proof}
Since $X$ is a groupoid, the maximal subgroupoid of $\Fun(X,\mathcal C)$
is canonically $\Map(X,\mathcal C^{\simeq})$.  Exponential adjunction
identifies the other displayed map with the map on maximal subgroupoids
of the arrow categories of the two functor categories.  For any
$\infty$-category $\mathcal A$ and objects $a,b$, its mapping space
$\Map_{\mathcal A}(a,b)$ is the homotopy fibre at $(a,b)$ of
\[
 \Fun(\Delta^1,\mathcal A)^{\simeq}\longrightarrow
 \mathcal A^{\simeq}\times\mathcal A^{\simeq}.
\]
An equivalence on the two displayed cores therefore gives essential
surjectivity and equivalences on every mapping space.  Conversely an
equivalence of categories gives both equivalences of cores.

For the last assertion, restriction commutes with cotensors.  Full
faithfulness on homotopy categories, applied to $M$ and $N^T$ for every
$\U$-small space $T$, identifies
$[T,\Map(M,N)]$ with $[T,\Map(f^*M,f^*N)]$.  Yoneda in the homotopy
category of spaces identifies the specified map of mapping spaces as an
equivalence.  Essential surjectivity on the homotopy category is the
same as essential surjectivity on objects.  All constructions are
canonical up to the contractible choices inherent in mapping objects.
For the adjoint formulation, an adjunction of $\infty$-categories
induces an adjunction of homotopy categories.  Full faithfulness
of a right adjoint is equivalent to invertibility of its counit,
and full faithfulness of a left adjoint to invertibility of its
unit.  Invertibility is detected in the homotopy category, so this
criterion upgrades full faithfulness directly; essential
surjectivity is unchanged.  This applies in particular to
restriction into $\mathcal P_R$, which has the Kan adjoints
although its entire mapping spaces need not be $\U$-small.
\end{proof}

\begin{claim}{E.2}{Proved}{An ordinary derivator realization for linear coefficients}
Fix universes $\U\subset\V\subset\mathcal W$, and an
$\U$-small $E_\infty$-ring spectrum $R$.  Let $\mathcal P_R$ denote the
$\infty$-category of $\U$-presentable $R$-linear stable
$\infty$-categories and $R$-linear colimit-preserving functors.  We regard
it as a category in the larger universes; all limits and colimits required
below have $\U$-small shapes.  Let
\[
 \D_R(A)=\Ho\Fun(NA^{\op},\Mod_R),\qquad
 \D'_R(A)=\Ho\Fun(NA^{\op},\mathcal P_R).
\]
Both are derivators on $\Cat_\U$, and their classes of functors satisfy
\[
                    W_{\D'_R}=W_{\D_R}^{\mathrm{lc}}.             \tag{E.2.1}
\]
Thus the equality in Part E has an ordinary-derivator realization for
commutative rings, by taking $R=Hk$.  It uses the specified symmetric
monoidal enhancement of $\D_R$.  It is functorial in maps of
$E_\infty$-rings by extension of scalars; it is not an assignment
constructed from an arbitrary unadorned derivator.
\end{claim}
\begin{proof}
We first check the size and existence assertions.  The category of
presentable categories admits small limits, computed on underlying
categories, and small colimits, computed by taking limits of right
adjoints \citep[Proposition 5.5.3.13 and Theorem 5.5.3.18]{EF_LurieHTT}.
It is symmetric monoidal, with tensor preserving colimits separately
\citep[Proposition 4.8.1.15 and Remark 4.8.1.24]{EF_LurieHA}.
We may identify $\mathcal P_R$ with
$\Mod_{\Mod_R}(\operatorname{Pr}^{L})$.  Limits and colimits of these
modules exist by \citet[Corollaries 4.2.3.3 and 4.2.3.5]{EF_LurieHA}.
The $\Mod_R$-action makes their underlying categories stable; the
corresponding assertion for spectra is
\citet[Proposition 4.8.2.18]{EF_LurieHA}.
The standard pointwise Kan-extension construction consequently gives
both displayed derivators.  We use $\mathcal W$ for their values and do
not assert that $\mathcal P_R$ itself is $\U$-presentable.
For a map $R\to S$ of $E_\infty$-rings, the coefficient functor is
\[
 \mathcal P_R\longrightarrow\mathcal P_S,
 \qquad \mathcal A\longmapsto
              \Mod_S\otimes_{\Mod_R}\mathcal A.
\]
Relative tensor products are characterized by balanced
colimit-preserving functors.  Their universal property supplies the
associative change-of-scalars comparisons for $R\to S\to T$ and
the identity comparisons coherently, so this gives functoriality in
the enhanced ring datum and induced morphisms of the displayed
derivators.  No choice of a ring presenting an arbitrary derivator
is implicit in this assertion.

Write $X=BA$ and $Y=BB$ for the source and target of the map of spaces
induced by a functor of indexing categories.  For $\mathcal A,\mathcal
B\in\mathcal P_R$, the constant-coefficient mapping space is
\[
 \Map\bigl(Y,\Fun_R^L(\mathcal A,\mathcal B)^{\simeq}\bigr).
                                                               \tag{E.2.2}
\]
Here $\Fun_R^L$ denotes the category of $R$-linear left adjoints and all
$R$-linear natural transformations.  Constant diagrams factor through
groupoid completion; this explains the use of $BB$ rather than $NB$ in
(E.2.2).  A fixed choice of inversion identifies the presheaf variance
with the covariant groupoid formula, naturally up to equivalence.

Suppose first that $f$ is in $W_{\D'_R}$.  For every $\mathcal B$ and
space $T$, its cotensor $\mathcal B^T=\Fun(T,\mathcal B)$ is again an
object of $\mathcal P_R$.  Applying the constant tests with
$\mathcal A=\Mod_R$ and this cotensor, and using
\[
 \Fun_R^L(\Mod_R,\mathcal B)\simeq\mathcal B,
                                                               \tag{E.2.3}
\]
we obtain bijections
$[T,\Map(Y,\mathcal B^{\simeq})]\to
[T,\Map(X,\mathcal B^{\simeq})]$ for every $T$.
Formula (E.2.3) is evaluation at $R$; its inverse sends $b$ to
$M\mapsto M\otimes_R b$.  It is the unit case of the linear functor
theorem \citep[Theorem 4.8.4.1]{EF_LurieHA}.
Yoneda gives equivalence on these mapping spaces.  Although the cores
may have a larger set of components, each component is $\U$-small:
its loop spaces are the equivalences in mapping spaces of a
locally $\U$-small presentable category, and the same is true after
a $\U$-small functor shape.  The test $T=*$ gives the component
bijection; Yoneda with connected $\U$-small $T$ then applies
component by component.  Taking successively
$\mathcal B=\Mod_R$ and $\mathcal B=\Fun(\Delta^1,\Mod_R)$, which is
presentable, stable, and $R$-linear, proves by E.1 that
$\Fun(Y,\Mod_R)\to\Fun(X,\Mod_R)$ is an equivalence.  This proves one
inclusion in (E.2.1).

For the converse we spell out the monodromy argument needed uniformly
for all coefficient categories.  For a pointed connected space $(Z,z)$
and a cocomplete $R$-linear category $\mathcal M$, evaluation
$U_z:\Fun(Z,\mathcal M)\to\mathcal M$ has a left adjoint $F_z$, given
by left Kan extension.  Its monad is
\[
 U_zF_z(m)=R[\Omega_zZ]\otimes_R m,
 \qquad R[\Omega_zZ]=R\otimes\Sigma^
 {\infty}_{+}\Omega_zZ.                                      \tag{E.2.4}
\]
Composition in the monad is path concatenation.  Evaluation is
conservative because $Z$ is connected, and preserves geometric
realizations because colimits of functors are pointwise.  The
$\infty$-categorical Barr--Beck theorem
\citep[Theorem 4.7.3.5]{EF_LurieHA} therefore identifies
$\Fun(Z,\mathcal M)$, over $\mathcal M$, with modules for the monad
in (E.2.4).  This proof also applies when $\mathcal M$ is cocomplete
without being presentable.  A current systematic treatment is
\citet[Proposition 1.2.7 and Corollary 1.2.8]{EF_PPS2025}.

If $R\ne0$, equivalence of $\Mod_R$-valued local systems forces a
bijection on components.  An omitted target component makes the
identity and zero endomorphisms of the constant $R$ on that component
indistinguishable after restriction.  Two source components lying over
one target component prevent the source system with fibres $0$ and $R$
on those components from being in the essential image.  These arguments
also handle empty source and target.
Choose a basepoint in each source component.  Restriction, being an
equivalence compatible with evaluation, identifies the left adjoints of
evaluation.  The canonical mate
$F_x\to f^*F_{f(x)}$ is consequently an equivalence.  Evaluating it at
$R$ and then at $x$ shows that the specified algebra map
\[
 R[\Omega_xX]\longrightarrow R[\Omega_{f(x)}Y]                 \tag{E.2.5}
\]
is an equivalence of $E_1$-algebras, not merely a Morita equivalence.
By (E.2.4), restriction on local systems in every cocomplete $R$-linear
$\mathcal M$ is therefore an equivalence on each component.

Apply this with $\mathcal M=\Fun_R^L(\mathcal A,\mathcal B)$.
Its colimits are computed pointwise: colimits commute with colimits and
with the $R$-action.  It is a cocomplete $R$-linear category.  Taking
maximal subgroupoids and then components proves invariance of every
constant-coefficient hom-set (E.2.2), as required.  If $R=0$, both
$\Mod_R$ and $\mathcal P_R$ are terminal categories, and both classes
contain every map.  This proves the remaining edge case.
\end{proof}

\begin{claim}{E.3}{Proved}{The two-categorical version of the linear realization}
Retain the hypotheses of E.2, and retain the full $\infty$-categories of
natural transformations between $R$-linear left adjoints.  The class
of maps $f:X\to Y$ for which all restrictions
\[
 \Fun\bigl(Y,\Fun_R^L(\mathcal A,\mathcal B)\bigr)
 \longrightarrow
 \Fun\bigl(X,\Fun_R^L(\mathcal A,\mathcal B)\bigr)               \tag{E.3.1}
\]
are equivalences is also $W_{\D_R}^{\mathrm{lc}}$.  The same is true
if equivalence in (E.3.1) is tested on ordinary homotopy categories.
These are the Hom-category constant tests in the coefficient
$(\infty,2)$-category $\mathcal P_R$.  No assertion that its diagram
assignment satisfies every axiom of Di Vittorio's 2-derivators is
needed for this equality.
\end{claim}
\begin{proof}
The forward implication in E.2 gave equivalence of (E.3.1), before
taking its core.  Conversely the single pair
$\mathcal A=\mathcal B=\Mod_R$ gives
$\Fun_R^L(\Mod_R,\Mod_R)\simeq\Mod_R$, so (E.3.1) already includes
the defining test for $W_{\D_R}^{\mathrm{lc}}$.
For the homotopy-category formulation the same converse is immediate,
and the forward implication follows from the equivalence already
proved.  Thus the ordinary realization and the Hom-category
realization agree in this case.  Passing to homotopy categories at one
shape loses information, but the cotensor and arrow tests in E.2
explain exactly why that observation does not obstruct the ordinary
realization.
\end{proof}

\begin{claim}{E.4}{Proved}{A specified second linear stage detects degree-two torsion}
For $R=H\Q$ use the particular coefficient category $\mathcal P_\Q$ of
E.2, and define
\[
 W_{\Q}^{(2)}=W_{\D'_\Q}^{\mathrm{lc}},\qquad
 W_{\Q}^{(1)}=W_{\D_\Q}^{\mathrm{lc}}=W_{\D'_\Q}.
\]
For every prime $p$, the projection $K(\Z/p,2)\to *$ is not in
$W_{\Q}^{(2)}$.  The projection $K(\Z,2)\to *$ is not in that
class either.  The first projection is in $W_{\Q}^{(1)}$, so the
inclusion $W_{\Q}^{(2)}\subset W_{\Q}^{(1)}$ is strict.
\end{claim}
\begin{proof}
For a field extension $L/\Q$, multiplication by a unit of $L$ is an
invertible $\Q$-linear natural transformation of the identity of
$\Mod_L$.  Scalar multiplication and composition give a coherent
map
\[
             B^2(L^\times)\longrightarrow\mathcal P_\Q^{\simeq}
                                                               \tag{E.4.1}
\]
based at $\Mod_L$: the single 1-morphism acts by the identity functor
and each 2-morphism acts by its scalar natural automorphism.  This may
be constructed with complexes of $L$-vector spaces, where all the
coherence equations are literal identities, and then localized.
On $\pi_2$, (E.4.1) sends a unit to its class as an automorphism of
the identity functor.  This homomorphism is injective: evaluating at
the free module $L$ sends it to that scalar in
$\pi_0\operatorname{Aut}_{\Mod_L}(L)=L^\times$.

For $A=\Z/p$ choose $L=\Q(\zeta_p)$ and the homomorphism
$A\to L^\times$ sending $1$ to a primitive $p$th root of unity.
Composing $K(A,2)=B^2A\to B^2L^\times$ with (E.4.1) gives a local
system of $\Q$-linear categories which is not constant: its induced
map on $\pi_2$ is nonzero, whereas every constant system has zero
map on $\pi_2$.  For $A=\Z$, use $L=\Q$ and $n\mapsto2^n$.
An equivalence to a constant system could change the identification
of its fibre by an autoequivalence, but conjugation acts by an
isomorphism on $\pi_2$ and cannot turn a nonzero homomorphism into
zero.  Thus the constant-system functor is not essentially surjective
in either case.

Finally $K(\Z/p,2)$ is simply connected and rationally acyclic.
One direct verification uses the path fibration with fibre
$K(\Z/p,1)$: transfer makes the latter rationally acyclic, and the
Serre spectral sequence for the contractible total space then gives
the assertion.  The rational universal-cover criterion of Part C
therefore puts its projection in $W_{\Q}^{(1)}$.
The inclusion of the two stages follows in general because constants
are locally constant.  The character mechanism above agrees with
the algebraically closed-field calculation of
\citet[Proposition 2.21]{EF_PPS2025}; our use of the coefficient
category $\Mod_{\Q(\zeta_p)}$ permits every prime over $\Q$.
\end{proof}

\begin{claim}{E.5}{Known}{The explicitly specified tower starting with sets}
For $r\ge0$ put $\mathcal C_r=\mathcal S_{\le r}$ and let
$\D_r$ be its derivator.  Then
\[
 W_{\D_r}=W_r,\qquad W_{\D_r}^{\mathrm{lc}}=W_{r+1},\qquad
 \bigcap_{r\ge0}W_r=W_\infty.                                 \tag{E.5.1}
\]
Consequently the stipulated choice $\mathcal C_r\mapsto
\mathcal C_{r+1}$ realizes the descending tower in Lead 4 for sets.
It does not characterize a unique coefficient-raising operation.
\end{claim}
\begin{proof}[Identification and intersection]
The first two identities are the truncation computations in Part C.
At $r=0$, the identification of the reflector with $\Pi_1$ is the
checked CATHOT C.1--C.4; we do not repeat that proof.  A map in all
$W_r$ is bijective on components and induces isomorphisms on every
homotopy group at each basepoint.  It is therefore a weak equivalence,
by the definition of weak equivalence of spaces.  The reverse inclusion
is immediate.  Empty spaces cause no exception because $W_0$ already
detects components.  CATHOT C.2 and B.11 also show why retaining only
categories modulo zigzags of natural transformations cannot produce
the subsequent stages: that representability theory stops at 1-types.
\end{proof}

\begin{claim}{E.6}{Known}{What the attached 2-derivator paper actually proves}
In arXiv:2309.05216v2, Di Vittorio's shapes are small simplicially
enriched categories, and a 2-prederivator is a trihomomorphism
$\mathrm{Dia}^{\op}\to\mathrm{GRAY}$, not a prederivator on
ordinary $\Cat$ alone.  For the 2-prederivator of a combinatorial
$\mathrm{sSet}_{\mathrm{Joyal}}$-enriched model category, axioms
(HDer 1)--(HDer 4) are proved, whereas (HDer 5) is Conjecture 3.2.23.
All five axioms, when assumed, are preserved by shift.
\end{claim}
\begin{proof}[Exact source check]
The shapes and 2-prederivators are Definitions 3.1.1 and 3.1.5; the
model construction is Theorem 3.1.14.  The four model verifications
are Propositions 3.2.3, 3.2.7, 3.2.11, and 3.2.18.
(HDer 1) concerns coproducts; (HDer 2) detection of equivalences of
1-cells; (HDer 3) left and right biadjoints; (HDer 4) exactness of the
weighted-collage squares.  (HDer 5) requires smothering underlying
diagram 2-functors for the walking arrow, walking isomorphism, and
walking adjunction.  The paper explicitly leaves the needed fullness
and conservativity statements conjectural, again in Section 4.2(2).
Shift invariance is Propositions 3.2.2, 3.2.6, 3.2.10, 3.2.16, and
3.2.22 \citep{EF_DiVittorio2026}.  The checked arXiv record identifies
v2 as the accepted version; we do not silently replace its conjecture
by an asserted theorem.

Our ordinary derivator $\D'_R$ satisfies ordinary Der 1--4 because
its enhanced coefficient category is complete and cocomplete.  E.3
concerns Hom-categories on groupoid shapes.  We have not constructed
an extension of it to Di Vittorio's enriched shapes, nor verified
his (HDer 5).  It would therefore be incorrect to call E.2 or E.3
a construction of a 2-derivator in his sense without additional work.
\end{proof}

\begin{claim}{E.7}{Proved}{A conditional realization by localized spaces}
Let $\mathcal L\subset\mathcal S$ be an accessible reflective
localization with reflector $L$, and let $W$ be the maps inverted by
$L$.  Assume $W$ is preserved by product with every small space.
For the derivator $\D_{\mathcal L}$ of $\mathcal L$, its
constant-coefficient class is $W$. In particular, if
$W_\D^{\mathrm{lc}}$ is accessible as a fundamental localizer,
there exists a derivator of localized spaces $\D^{\mathrm{loc}}$
with
\[
                 W_{\D^{\mathrm{loc}}}=W_\D^{\mathrm{lc}}.
\]
The construction is functorial for specified inclusions between
accessible fundamental localizers. By B.4, the accessibility
hypothesis holds whenever $\D$ is the derivator of a
$\U$-presentable infinity-category.
\end{claim}
\begin{proof}
For a local space $B$ and any space $A$, the space $\Map(A,B)$ is
local: testing a map $f\in W$ against it is the same as testing
$f\times A$ against $B$.  Thus each mapping space between local
objects is local, and a map in $W$ preserves all the required
constant-coefficient hom-sets.  Conversely the terminal space is
local.  Testing constants $*$ and all $B\in\mathcal L$ gives
bijections
$\Hom_{\Ho\mathcal L}(LY,B)\to
\Hom_{\Ho\mathcal L}(LX,B)$.
Yoneda in $\Ho\mathcal L$ says that the specified $Lf$ is an
isomorphism there, and hence an equivalence.  This is precisely
$f\in W$.  The accessible localization is presentable and complete,
so its associated derivator exists.

For the additional assertion, B.3 supplies the fundamental-localizer
axioms for $W_\D^{\mathrm{lc}}$. Every fundamental localizer is
stable under products: for $f:A\to B$ in the localizer, apply the
relative Theorem A axiom to $f\times C$ over $C$. Its slices are
$f\times(C/c)$, and their projections to $f$ are ordinary weak
homotopy equivalences since $C/c$ has a terminal object.
Two-out-of-three gives the required slice equivalences. The
correspondence for the test category $\Delta$ in
\citet[Theorems~4.2.15 and~1.4.3]{Cisinski2006} sends an accessible fundamental
localizer to an accessible localization of spaces with exactly the
same specified weak equivalences under categorical realization.
Its class is product stable by the preceding argument. The first
part of this proof therefore applies. If $W\subseteq W'$ are
specified accessible fundamental localizers, restriction of the
$W'$ reflector to $W$-local spaces is the canonical colimit-preserving
functor $\mathcal S_W\to\mathcal S_{W'}$. Uniqueness in the
universal property of localization supplies coherent composition.
This is the asserted functoriality on the poset of localizers;
an arbitrary morphism of derivators need not supply such an inclusion.
\end{proof}

\begin{claim}{E.8}{Proved}{The complete second-stage loop criterion}
Let $R\ne0$ be an $E_\infty$-ring spectrum, and put
$W_R^{(2)}=W_{\D'_R}^{\mathrm{lc}}$ for the realization in E.2.
A map $f:X\to Y$ of $\U$-small spaces is in $W_R^{(2)}$ if
and only if it is bijective on components and, for one basepoint
$x$ in every component, the specified map
\[
 \Omega_xf:\Omega_xX\longrightarrow\Omega_{f(x)}Y
\]
belongs to $W_{\D_R}^{\mathrm{lc}}$.  In particular, for a
nonzero commutative ordinary ring $k$,
\[
 f\in W_k^{(2)}
 \quad\Longleftrightarrow\quad
 \begin{cases}
 \pi_0(f)\text{ is bijective},\\
 \pi_1(f),\pi_2(f)\text{ are isomorphisms at every basepoint},\\
 H_*(X\langle2\rangle;k)\longrightarrow
 H_*(Y\langle2\rangle;k)\text{ is an isomorphism on each component}.
 \end{cases}                                                  \tag{E.8.1}
\]
Here $X\langle2\rangle$ is the homotopy fibre at the chosen
basepoint of $X\to\tau_{\le2}X$; all maps of covers are the
ones induced by $f$ and the chosen basepoints.  This is a
full computation of the specified rational second stage on
taking $k=\Q$, with no nilpotence or finite-type assumption on
$X$ or $Y$.  In that case the last condition is equivalently
that $\pi_j(f)\otimes\Q$ is an isomorphism for every $j\ge3$;
the 2-connected covers to which rational Whitehead applies are
automatically simply connected.
\end{claim}
\begin{proof}
First suppose $X,Y$ are connected and pointed.  The symmetric
monoidal functor
\[
 \mathcal S_\U\longrightarrow\mathcal P_R,
 \qquad T\longmapsto\mathcal P_R[T]
       :=\colim_{t\in T}\Mod_R
\]
sends Cartesian products to relative tensor products, by the
universal property of colimits and the separate preservation of
colimits by the tensor product.  Its unit is $\Mod_R$.
The free $R$-linear presheaf universal property identifies
\[
 \mathcal P_R[T]\simeq\Fun(T^{\op},\Mod_R).                   \tag{E.8.2}
\]
To check this identification, an $R$-linear left adjoint from
either side to $\mathcal A$ is precisely a functor
$T\to\mathcal A$: on the right it is the colimit-preserving
extension of its values on the $R$-linear representables; on
the left it is the compatible family of functors
$\Mod_R\to\mathcal A$ determined by those same values.
These universal properties identify both functors and natural
transformations.  For a map $g:T\to T'$, the induced map in
(E.8.2) is left Kan extension $g_!$; its right adjoint is the
specified restriction $g^*$.  Thus the comparison is natural
in the direction required here, not just an abstract equivalence
of categories.  For groupoids the inversion equivalence may
remove the opposite.

Evaluation $U_x:\Fun(X,\mathcal P_R)\to\mathcal P_R$ is
conservative and preserves geometric realizations.  It has a
left adjoint $F_x$ by pointwise Kan extension.  Barr--Beck
therefore applies even though $\mathcal P_R$ need not itself
be presentable in the indexing universe.  The free-evaluation
monad is
\[
 U_xF_x(\mathcal A)
       =\mathcal P_R[\Omega_xX]\otimes_{\Mod_R}\mathcal A.
                                                               \tag{E.8.3}
\]
The $E_1$-algebra structure of $\mathcal P_R[\Omega_xX]$
comes from loop multiplication, and (E.8.3) is compatible with
the full monad structure.  For $f$ the canonical comparison
of these monads is tensoring with the specified algebra map
\[
 a_f:\mathcal P_R[\Omega_xX]\longrightarrow
                     \mathcal P_R[\Omega_{f(x)}Y].            \tag{E.8.4}
\]
If restriction of $\mathcal P_R$-valued local systems is an
equivalence, its compatibility with evaluation identifies the
left adjoints of evaluation, so the comparison of monads is an
equivalence.  Evaluate it at the unit $\Mod_R$ to see that
(E.8.4) is an equivalence of underlying categories, and hence
of $E_1$-algebra objects.  Conversely equivalence of (E.8.4)
makes the monads equivalent, and the monadic comparison
identifies their restriction functor with $f^*$; consequently
that specified restriction is an equivalence.  E.1's adjunction
argument identifies this with the ordinary-derivator definition.

By (E.8.2), (E.8.4) is left Kan extension along $\Omega_xf$.
A left adjoint is an equivalence exactly when its right
adjoint is, so (E.8.4) is an equivalence exactly when
$\Omega_xf$ belongs to $W_{\D_R}^{\mathrm{lc}}$.
For disconnected spaces the same zero-versus-unit argument
as in E.2 forces a component bijection.  All functor categories
then decompose as products over the corresponding components,
and the connected proof applies to each factor.  It also
shows that an empty source is permitted precisely when the
target is empty.

For $R=Hk$, apply the loop-algebra computation of Part C to
$\Omega_xf$.  Its component bijection is exactly the
isomorphism on $\pi_1(X,x)$, and its fundamental-group
isomorphism is the isomorphism on $\pi_2(X,x)$.
The universal cover of the identity component of
$\Omega_xX$ is $\Omega(X\langle2\rangle)$: loop the
fibre sequence $X\langle2\rangle\to X\to\tau_{\le2}X$,
then take the identity component.  Other components are
identified by translations, and $\Omega_xf$ commutes with
these translations coherently.  Finally, on simply connected
spaces, a specified map is a $k$-homology equivalence if and
only if its based loop map is a $k$-homology equivalence, as
proved by the connected bar--cobar comparison in Part C.
Apply this to the 2-connected covers to obtain (E.8.1).
\end{proof}

\begin{claim}{E.9}{Proved}{Degree-three torsion survives the second rational stage}
For every prime $p$, the projection $K(\Z/p,3)\to *$ belongs to the
specific $W_\Q^{(2)}$ of E.4.  More generally this holds for every
nonempty, 2-connected, rationally acyclic space in place of
$K(\Z/p,3)$.  In particular
\[
             W_\infty\subsetneq W_\Q^{(2)}
                  \subsetneq W_\Q^{(1)}.
\]
\end{claim}
\begin{proof}
Let $\mathcal A\in\mathcal P_\Q$.  For $i\ge3$ the homotopy
groups of $\mathcal P_\Q^{\simeq}$ at $\mathcal A$ are
\[
 \pi_i(\mathcal P_\Q^{\simeq},\mathcal A)
 \cong \pi_{i-2}
   \operatorname{End}_{\Fun_\Q^L(\mathcal A,\mathcal A)}
                         (\id_{\mathcal A}),                 \tag{E.9.1}
\]
where $\operatorname{End}$ denotes the mapping spectrum.  Indeed the
loop space at $\mathcal A$ is the space of its autoequivalences; the
loop space there at the identity is the space of invertible natural
endomorphisms of that identity.  In positive homotopy degrees this is
the same as the underlying space of its endomorphism spectrum.  The
functor category is $\Q$-linear and stable, so the groups in
(E.9.1) are $\Q$-vector spaces.

The endpoint map
\[
 \Fun(\Delta^1,\mathcal P_\Q)^{\simeq}
       \longrightarrow (\mathcal P_\Q^{\simeq})^2
\]
has homotopy fibre
$\Fun_\Q^L(\mathcal A,\mathcal B)^{\simeq}$ at
$(\mathcal A,\mathcal B)$.  Its homotopy groups in degrees at least
two are homotopy groups of rational mapping spectra, hence are
$\Q$-vector spaces.  The long exact homotopy sequence shows that
the total space has $\Q$-vector-space homotopy groups in every
degree at least three as well: kernels and cokernels of maps between
uniquely divisible abelian groups are uniquely divisible, and an
extension of two such groups is uniquely divisible.

We use the following elementary Postnikov observation.  If $T$ is
2-connected and rationally acyclic, and $K$ is any space with
$\pi_i(K,k)$ a $\Q$-vector space for $i\ge3$, then evaluation
$\Map(T,K)\to K$ is an equivalence.  Its fibre at $k$ is
$\Map_*(T,K)$, which agrees with the pointed mapping space into
the 2-connected cover of $(K,k)$, since
$\Map_*(T,\tau_{\le2}K)$ is contractible.  For a $\Q$-vector
space $V$ all reduced cohomology groups of $T$ with coefficients
in $V$ vanish.  Hence $\Map_*(T,K(V,n))$ is contractible for
every $n$.  Induction over the Postnikov fibrations proves the same
for every finite Postnikov section of the 2-connected cover; taking
their inverse limit proves the observation.  This argument uses
Postnikov completeness of spaces, and does not replace a limit by a
colimit or assume finite type.  If $K$ is empty both sides are empty;
otherwise the argument applies independently to all components.

Apply the observation to the two spaces $K_0,K_1$ in E.1 for
$\mathcal C=\mathcal P_\Q$.  It proves that the constant functor
$\mathcal P_\Q\to\Fun(T,\mathcal P_\Q)$ is an equivalence.
Finally $K(\Z/p,3)$ is 2-connected and rationally acyclic.  Transfer
gives rational acyclicity of $K(\Z/p,1)$, and two applications of
the path-fibration Serre spectral sequence give it for degrees two
and three.  Its nonzero $\pi_3$ shows that its projection is not a
weak equivalence.  Strictness of the other inclusion is E.4.
\end{proof}

\begin{claim}{E.10}{Proved}{A rational loop tower and its ordinary coefficient realization}
Define the following specified tower of classes of maps of spaces:
\[
 V_0=W_{\D_\Q},\qquad
 V_{r+1}=\{f:X\to Y\mid \pi_0(f)\text{ is bijective and }
                \Omega_xf\in V_r\text{ for every component}\}.
                                                               \tag{E.10.1}
\]
For $r\ge1$ its exact description is
\[
 f\in V_r\quad\Longleftrightarrow\quad
 \begin{cases}
 \pi_0(f)\text{ is bijective},\\
 \pi_i(f)\text{ is an isomorphism }(1\le i\le r),\\
 \pi_j(f)\otimes\Q\text{ is an isomorphism }(j>r).
 \end{cases}                                                \tag{E.10.2}
\]
The rationalization condition concerns the higher abelian homotopy
groups and their specified maps.  We have
$V_1=W_\Q^{(1)}$, $V_2=W_\Q^{(2)}$ for the actual coefficient
categories constructed above, strict inclusions
$V_{r+1}\subsetneq V_r$ for every $r\ge0$, and
$\bigcap_{r\ge0}V_r=W_\infty$.
Every $V_r$ is an accessible fundamental localizer. For $r\ge1$
let $\mathcal C_r=\mathcal S_{V_r}$ be its category of local spaces,
and write $\D_r=\D_{\mathcal C_r}$, with $\D_0=\D_\Q$.
These are a specified tower of ordinary derivators satisfying
\[
                  W_{\D_r}=V_r,
       \qquad W_{\D_r}^{\mathrm{lc}}=V_{r+1}
       \quad(r\ge0).
 \tag{E.10.3}
\]
One may instead use $\D_1=\D'_\Q$ from E.2 and retain
$\D_r=\D_{\mathcal C_r}$ for $r\ge2$; the same equalities hold.
Thus the tower, including its intersection, has a construction by
ordinary coefficient categories. The stages $\mathcal C_r$ are
localized spaces; they are not identified here with categories of
higher $\Q$-linear categories.
More generally, if an initial derivator $\mathbb E$ has accessible
$W_{\mathbb E}^{\mathrm{lc}}\subseteq W_0$, the same repeated
localized-space realization constructs an ordinary coefficient
tower whose intersection is $W_\infty$. In particular this applies
to every nonterminal $\U$-presentable infinity-category coefficient.
\end{claim}
\begin{proof}
For a group-like loop space, rational homology equivalence detects
its component group by $H_0$, and translation identifies its
components.  The loop-chain criterion in Part C consequently
identifies $V_1$ with $W_{\D_\Q}^{\mathrm{lc}}$.  Rational
Whitehead on the simply connected covers gives (E.10.2) at
$r=1$.  E.8 identifies the next loop step with the actual
second categorical stage.  In general
$\pi_i(\Omega_xX)=\pi_{i+1}(X,x)$ for $i\ge1$, while
$\pi_0(\Omega_xX)=\pi_1(X,x)$ as a group.  Apply the inductive
criterion to all components of the loop map; translations make
this equivalent to the criterion at its identity component.
This proves (E.10.2) for $r+1$ and proves all the inclusions.
The inclusion $V_1\subset V_0$ also follows from the
constant-versus-locally-constant inclusion.

For any prime $p$, the projection $K(\Z/p,r+1)\to *$ is in
$V_r$ but not $V_{r+1}$.  For $r=0$ use finite-group transfer;
for $r\ge1$ use (E.10.2), since its only nonzero homotopy
group is the torsion group in degree $r+1$.  Thus every inclusion
is strict.  A map in every $V_r$ is bijective on components and
is an isomorphism on each homotopy group by taking $r$ large
enough for that degree.  It is therefore a weak equivalence.
The reverse inclusion is immediate from (E.10.1).

We now construct (E.10.3). We first record the coefficient calculation
which makes the induction possible. Let
$L:\mathcal S\to\mathcal C\subseteq\mathcal S$ be an accessible
reflective localization whose inverted maps $W$ are product stable,
and suppose $W\subseteq W_0$. The point is always local, and
the empty space is local because the localization unit belongs to
$W$ and therefore preserves its empty set of components. These
two local spaces are distinct. For any connected pointed
space $(X,x)$, evaluation
\[
 U_x:\Fun(X,\mathcal C)\longrightarrow\mathcal C
\]
has a left adjoint, is conservative, and preserves geometric
realizations. Its monad is
\[
                  M\longmapsto L(\Omega_xX\times M).
 \tag{E.10.4}
\]
Indeed the free local system has value the copower by the path
space, and colimits of local spaces are obtained by applying $L$
to the corresponding colimits of spaces. Barr--Beck applies as in
E.8. For a pointed map $f:X\to Y$ between connected spaces, the canonical comparison of
the monads (E.10.4) is an equivalence if and only if
$\Omega_xf\in W$. Necessity follows by evaluating the comparison
at $M=*$. For sufficiency, product stability makes
$\Omega_xf\times M$ a member of $W$ for every local $M$,
so the comparison is an equivalence on every object, compatibly
with its monad structure. Monadicity identifies its restriction
functor with the specified $f^*$, exactly as in E.8.

For arbitrary spaces, an equivalence of local-system categories
forces a component bijection. If a target component is missed,
restriction cannot detect the difference between the empty and
point-valued constant local systems on that component. If two
source components map into one target component, the local system
which has empty value on one and point value on the other cannot
be a restriction: the fibers of a local system on a connected
space are equivalent. Once components correspond, apply the
connected calculation componentwise. This proves
\[
 W_{\D_{\mathcal C}}^{\mathrm{lc}}
   =\{f:\pi_0(f)\text{ is bijective and }
                   \Omega_x f\in W\text{ on every component}\}.
 \tag{E.10.5}
\]
E.1 identifies the enhanced equivalence used in this argument with
the ordinary homotopy-category definition.

The class $V_1=W_{\D_\Q}^{\mathrm{lc}}$ is an accessible
fundamental localizer by B.3--B.4. Apply E.7 to form
$\mathcal C_1$. It is presentable, its constant class is $V_1$,
and (E.10.5) identifies its locally constant class with $V_2$.
B.3--B.4 now make $V_2$ accessible and fundamental, so E.7 forms
$\mathcal C_2$. Inductively, the same argument constructs every
$\mathcal C_r$ and proves (E.10.3). The base $V_0$ is the
constant-coefficient class of $\D_\Q$ and is accessible by B.4.
The component condition needed in (E.10.5) follows at every stage
from (E.10.2). Finally E.2 and E.8 say respectively
$W_{\D'_\Q}=V_1$ and $W_{\D'_\Q}^{\mathrm{lc}}=V_2$,
which proves the assertion about the alternative first stage.

For the final assertion start with
$U_1=W_{\mathbb E}^{\mathrm{lc}}$ and set
$U_{r+1}=\{f:\pi_0(f)\text{ bijective},\ \Omega_xf\in U_r\}$
for $r\ge1$. The same induction using (E.10.5) constructs the
coefficient categories and preserves accessibility. Since
$U_1\subseteq W_0$, induction on $r$ gives
$U_r\subseteq W_{r-1}$, where $W_{r-1}$ is equivalence on
$(r-1)$-types. Hence a map in every $U_r$ is an isomorphism on
every homotopy group and on components. Conversely every weak
homotopy equivalence belongs to every stage by A.2, so the
intersection is $W_\infty$. B.4 supplies accessibility for a
presentable initial coefficient, and G.1 supplies component
detection when that coefficient is nonterminal.
\end{proof}

\begin{claim}{E.11}{Open}{Arbitrary bare-derivator functoriality and higher linear coefficients}
For an arbitrary derivator on $\Cat_\U$ we have not constructed a
functorial assignment $\D\mapsto\D'$ with
$W_{\D'}=W_\D^{\mathrm{lc}}$.  A phrase such as ``$\D$-linear
categories'' does not specify such an assignment without a monoidal
enhancement and a choice of morphisms.  E.2 settles a specified
enhanced linear case, and B.4 together with E.7 gives a realizing
derivator of localized spaces for every presentable
infinity-category coefficient. Neither proves a functorial
assignment on all bare derivators and all their morphisms.
The full second linear stage is computed in E.8, and E.10 constructs
an entire ordinary coefficient tower with its stated intersection.
A realization of its further stages by higher $\Q$-linear
categories, with a specified higher monodromy construction and
morphism class, remains unresolved here.
\end{claim}
\begin{proof}[Remaining mathematical task]
The unenhanced problem asks for a construction on derivators and their
morphisms, not merely for a representative of an equivalence class of
localizers.  The failed step in extending E.2 is that neither the
symmetric monoidal category of modules nor its category of linear
categories is part of the data of an arbitrary derivator.
The next linear-categorical task is to construct an operation on
$\Q$-linear categories whose third weak-equivalence class can be compared with maps
inducing an equivalence on 3-types and rational homology
equivalences of 3-connected covers.  The ordinary category
$\mathcal P_\Q$ is not a presentable coefficient category in
its original universe, so repeating the notation
``modules in presentable categories'' is not by itself a
construction of that higher linear categorical level. E.10's
localized-space coefficients already give an ordinary realization
of the same third class and of every subsequent loop class.
The higher monodromy
equivalences of \citet[Theorem 3.2.24]{EF_PPS2025} provide a precise
framework, but their general higher-Hochschild statement is conditional
on their Conjecture 3.3.1.  We do not import that conditional statement
as a theorem or assert that arbitrary iteration has been defined.
\end{proof}
